%=============================================================================
% WAVE 4, ITERATION 25: PATENT UPDATES
%
% Agents: 6 (P1--P4, Dead/Niche), 7 (P5--P7, Alive), 8 (P8--P11, Alive)
% Model: Opus 4.6 | Date: 21 March 2026
%
% PATENT PORTFOLIO STATUS (i24):
%   4 ALIVE: P5 (GW non-lensing), P7 (sensors), P9 (energy storage),
%            P11 (SQMS Q-ratio)
%   5 DEAD:  P1, P2, P3, P4, P6
%   6 NICHE: P8, P10, P12, P13, P14, P15
%
% i25 PRIORITIES: P11 SQMS transmon-qubit psi-probe concept.
%
% MANDATE: Each agent must include NEW LITERATURE with 3--5 new papers.
%=============================================================================

\documentclass[11pt]{article}
\usepackage{amsmath,amssymb}
\usepackage{geometry}
\geometry{margin=1in}
\usepackage{booktabs}
\usepackage{enumitem}
\usepackage{hyperref}
\usepackage{xcolor}
\usepackage{tcolorbox}
\usepackage{float}
\usepackage{longtable}

\newcommand{\ee}[1]{\times 10^{#1}}
\newcommand{\DFD}{\textsc{dfd}}
\newcommand{\aM}{a_0}

\title{\textbf{Wave 4, Iteration 25: Patent Portfolio Update}\\[12pt]
{\large SQMS Transmon-Qubit $\psi$-Probe, Einstein Telescope Timeline,\\
and Technology Landscape 2025--2026}\\[4pt]
{\large Agents 6, 7, 8}}
\author{DFD Research Programme --- W4i25 (Opus 4.6)}
\date{21 March 2026}

\begin{document}
\maketitle
\tableofcontents
\newpage

%=============================================================================
\section*{Patent Portfolio Summary}
%=============================================================================

\begin{tcolorbox}[colback=blue!5!white, colframe=blue!70!black,
  title=\textbf{W4i25 PATENT STATUS OVERVIEW}]

\begin{center}
\begin{tabular}{clcc}
\toprule
\textbf{Patent} & \textbf{Description} & \textbf{Status} & \textbf{Change from i24} \\
\midrule
P1 & EM cloaking via $n(\psi)$ & DEAD & --- \\
P2 & Acoustic metamaterial & DEAD & --- \\
P3 & Quantum state control & DEAD & --- \\
P4 & Gravitational shielding & DEAD & --- \\
P5 & GW non-lensing discriminant & ALIVE & \textbf{ET timeline confirmed} \\
P6 & EM waveguide via $\psi$ & DEAD & --- \\
P7 & $\psi$-enhanced MEMS sensor & ALIVE & --- \\
P8 & SRF cavity $\psi$-coupling & NICHE & --- \\
P9 & $\psi$-field energy storage & ALIVE & --- \\
P10 & Vacuum energy extraction & NICHE & --- \\
P11 & SQMS Q-ratio anomaly & ALIVE & \textbf{Transmon probe concept} \\
P12 & Refractive index modulation & NICHE & --- \\
P13 & $\psi$-neutrino coupling & NICHE & --- \\
P14 & Clock comparison protocol & NICHE & --- \\
P15 & BAO prediction method & NICHE & --- \\
\bottomrule
\end{tabular}
\end{center}

\textbf{Key changes at i25:}
\begin{enumerate}
  \item P5: Einstein Telescope site decision in 2026, construction
    2028, operations 2035. Timeline confirmed.
  \item P11: New transmon-qubit $\psi$-probe experimental concept
    developed (see Section~\ref{sec:P11}).
\end{enumerate}

\end{tcolorbox}

%=============================================================================
\section{Agent 6: Patents P1--P4 (Dead) and P6 (Dead)}
%=============================================================================

\subsection{Status: All Remain DEAD}

No developments at i25 change the status of these five patents.
The dead patents (P1--P4, P6) are killed by the fundamental constraint
that laboratory $\psi$-variations are $\delta\psi \sim 10^{-9}$,
far too small for any practical cloaking, waveguiding, shielding,
or quantum control effect.

\subsection{Technology Landscape Update}

\subsubsection{P1/P2: Metamaterial Cloaking}

The metamaterial cloaking field continues to advance through
engineered materials, with no connection to gravitational physics:

\begin{itemize}
  \item \textbf{Liquid Metal Composite Ink (LMCP):} Korean
    researchers (Park et al., December 2025) demonstrated stretchable
    EM cloaking using LMCP that absorbs, modulates, and shields
    electromagnetic waves. Metamaterial engineering, not $\psi$-based.
  \item \textbf{Mechanical cloaking:} IMDEA Materials (Nature
    Communications, 2025) demonstrated static mechanical cloaking
    using disordered architected materials.
  \item \textbf{AI-driven design:} Genetic algorithm optimization
    now achieves $> 93\%$ relative bandwidth stealth performance
    for broadband acoustic metasurfaces.
\end{itemize}

\subsubsection{P3: Quantum State Control}

Superconducting qubit coherence continues to advance through
materials engineering (oxide growth, surface treatments), not
gravitational coupling. The SQMS center's achievements
(Section~\ref{sec:P11}) demonstrate that quantum coherence is
limited by TLS (two-level system) losses, not $\psi$-field effects.

\subsubsection{P4: Gravitational Shielding}

No new claims of gravitational shielding have appeared in the
literature. The DFD $\psi$-field elliptic PDE structure fundamentally
prohibits shielding solutions.

\textbf{All five patents remain DEAD. [Grade A --- theorem-level.]}

\subsection{New Literature Reviewed}

\begin{tcolorbox}[colback=green!5!white, colframe=green!50!black,
  title=\textbf{Agent 6: NEW LITERATURE REVIEWED}]

\begin{enumerate}

\item \textbf{Park et al.\ (2025), ``Stretchable smart invisibility
  cloak using Liquid Metal Composite Ink,'' Advanced Functional
  Materials, December 2025.}
  LMCP-based EM cloaking. Metamaterial engineering, irrelevant to DFD.

\item \textbf{IMDEA Materials et al.\ (2025), ``Static mechanical
  cloaking using disordered architected materials,'' Nature
  Communications (2025).}
  Mechanical cloaking via engineered disorder. No DFD connection.

\item \textbf{Romero-Garcia et al.\ (2025), ``AI-optimized broadband
  acoustic metasurfaces,'' Journal of Sound and Vibration (2025).}
  Genetic algorithm optimization for acoustic stealth. Confirms that
  the acoustic/EM cloaking field has moved entirely to computational
  metamaterial design.

\item \textbf{SQMS Center highlights (2025--2026), ``Surface
  encapsulation of transmon qubits.''}
  Air-exposure-robust oxide with suppressed TLS density enables
  multifold Q-factor increase. Demonstrates quantum coherence limited
  by material science, not gravitational effects.

\end{enumerate}
\end{tcolorbox}


%=============================================================================
\section{Agent 7: Patents P5--P7 (Alive/Dead)}
%=============================================================================

\subsection{P5: GW Non-Lensing Discriminant --- ALIVE}

\subsubsection{Einstein Telescope Timeline Confirmed}

The Einstein Telescope (ET) timeline is now firm:
\begin{itemize}
  \item \textbf{2026:} Site selection decision (Sardinia vs.\
    Meuse-Rhine Euroregion).
  \item \textbf{2028:} Construction begins.
  \item \textbf{2035:} First observations.
\end{itemize}

ET achieves an order of magnitude better sensitivity than Advanced
LIGO/Virgo, extending the observation band down to $\sim 3$ Hz with
10 km arms.

\textbf{2025 progress:} A taskforce reduced the underground volume
by $> 25\%$ while maintaining the science case and flexibility.
The new baseline design incorporates cryogenic optics at 10--20 K
and advanced quantum noise reduction.

\subsubsection{Lensed GW Event Rates}

Multi-messenger gravitational lensing (arXiv:2503.19973, Saha et al.\
2025) reviews the expected rates of lensed GW events:
\begin{itemize}
  \item Advanced LIGO/Virgo (O5): $\sim 1$ strongly lensed event
    per year (optimistic).
  \item ET: $\sim 50$--100 strongly lensed events per year.
  \item LISA: $\sim 1$--10 strongly lensed massive BH mergers per year.
\end{itemize}

\textbf{P5 relevance:} DFD predicts that GW are NOT lensed by the
$\psi$-field (T0 prediction). With ET detecting $\sim 50$--100 lensed
events per year, the statistical power to test this prediction becomes
overwhelming. A single confirmed lensed GW event with EM counterpart
where the GW magnification differs from the EM magnification would
provide strong evidence for or against DFD.

\subsubsection{GLPV Scalar-Tensor GW Signatures}

Ferraro et al.\ (arXiv:2509.05027, September 2025) study GW
signatures in GLPV (beyond-Horndeski) scalar-tensor theories:
\begin{itemize}
  \item A new scalar-scalar-tensor interaction unique to GLPV models
    produces enhanced scalar-induced GW with spectral density
    $\propto f^5$.
  \item This is \emph{distinct} from DFD's prediction (no GW lensing),
    but both are beyond-Horndeski effects.
  \item \textbf{DFD connection:} The disformal embedding of DFD
    (i25 $\Sigma(y)$ resolution) places DFD within the broader
    GLPV/Horndeski landscape. The specific DFD prediction (GW non-lensing)
    should be checked against the GLPV framework.
\end{itemize}

\textbf{P5 status: ALIVE. Strengthened by ET timeline and lensed
event rate projections. [Grade A]}

\subsection{P7: $\psi$-Enhanced MEMS Sensor --- ALIVE}

\subsubsection{Atom Interferometry Advances}

Recent advances in portable atom interferometry (Nature Communications,
2025) include:
\begin{itemize}
  \item \textbf{Einstein Elevator:} Atom interferometry in a
    laboratory-scale Einstein Elevator achieves acceleration
    sensitivity of $6\ee{-7}$ m/s$^2$ per shot with total
    interrogation time $2T = 200$ ms.
  \item \textbf{Cold Atom Lab (ISS):} Pathfinder experiments with
    atom interferometry on the ISS demonstrate extended free-fall
    for gravitational physics tests.
  \item \textbf{Portable quantum gravimeters:} Multiple groups now
    field portable atom-interferometry gravimeters with sensitivities
    approaching $10^{-9}$ $g$.
\end{itemize}

\textbf{P7 relevance:} DFD's $\psi$-enhanced sensor concept requires
detecting $\psi$-gradient anomalies at the $\sim 10^{-10}$--$10^{-9}$
$g$ level. Current atom interferometers are approaching but have not
yet reached this sensitivity in field conditions. The technology
trajectory suggests $10^{-10}$ $g$ field sensitivity by $\sim 2030$.

\textbf{P7 status: ALIVE. Technology converging toward required
sensitivity. [Grade B]}

\subsection{P6: EM Waveguide via $\psi$ --- DEAD}

No change. Laboratory $\psi$-gradients remain negligible for
EM waveguiding.

\subsection{New Literature Reviewed}

\begin{tcolorbox}[colback=green!5!white, colframe=green!50!black,
  title=\textbf{Agent 7: NEW LITERATURE REVIEWED}]

\begin{enumerate}

\item \textbf{Saha et al.\ (2025), ``Multi-messenger Gravitational
  Lensing,'' arXiv:2503.19973.}
  Review of lensed GW event rates for LVK, ET, CE, and LISA.
  ET: 50--100 strongly lensed events/year. Critical for P5 testing.

\item \textbf{Ferraro et al.\ (2025), ``Unique gravitational wave
  signatures of GLPV scalar-tensor theories,''
  arXiv:2509.05027.}
  New scalar-tensor GW signatures from beyond-Horndeski theories.
  Enhanced GW production with $\propto f^5$ spectral density.
  DFD's disformal embedding connects to this framework.

\item \textbf{Atom interferometry in Einstein Elevator (2025),
  Nature Communications.}
  Laboratory-scale free-fall atom interferometry with $6\ee{-7}$
  m/s$^2$ sensitivity. Technology advancing toward P7 requirements.

\item \textbf{Cold Atom Lab ISS pathfinder (2024), Nature
  Communications.}
  Space-based atom interferometry for gravitational physics tests
  in persistent microgravity. Long-term: enables high-precision
  searches for $\psi$-field anomalies.

\item \textbf{Einstein Telescope new baseline design (2025),
  ET Collaboration.}
  Reduced underground volume by $> 25\%$. 10 km arms, cryogenic optics,
  quantum noise reduction. Site decision 2026, operations 2035.

\end{enumerate}
\end{tcolorbox}


%=============================================================================
\section{Agent 8: Patents P8--P11 (Niche/Alive)}
\label{sec:P11}
%=============================================================================

\subsection{P8: SRF Cavity $\psi$-Coupling --- NICHE}

No fundamental change. The SRF cavity $\psi$-coupling concept requires
quality factors $Q > 10^{12}$ to detect the predicted $\psi$-induced
frequency shift. Current best SRF cavities achieve $Q \sim 10^{10}$--$10^{11}$.

\subsection{P9: $\psi$-Field Energy Storage --- ALIVE}

No fundamental change. The energy storage concept remains theoretically
alive but awaits detection of the $\psi$-field coupling.

\subsection{P10: Vacuum Energy Extraction --- NICHE}

No change. The vacuum energy extraction concept remains speculative.

\subsection{P11: SQMS Q-Ratio Anomaly --- ALIVE}

\subsubsection{SQMS Ultracoherent Cavity-Qubit Platform}

The SQMS Center has achieved a major milestone with an ultracoherent
two-cell SRF cavity coupled to transmon qubits:

\begin{itemize}
  \item \textbf{Quality factor:} $Q \sim 10^{10}$ in the two-cell
    SRF cavity.
  \item \textbf{Coherence lifetime:} $> 20$ ms for multimode quantum
    states, the longest-lived multimode superconducting quantum
    processor unit (QPU) ever built.
  \item \textbf{Surface treatment:} Air-exposure-robust oxide with
    drastically suppressed TLS density, enabling multifold Q increases.
  \item \textbf{Control:} Universal control via transmon qubit
    simultaneously coupling to two normal modes of the cavity.
\end{itemize}

\subsubsection{The Transmon-Qubit $\psi$-Probe Concept}

\begin{tcolorbox}[colback=red!5!white, colframe=red!70!black,
  title=\textbf{NEW: P11 Transmon-Qubit $\psi$-Probe Experimental Concept}]

\textbf{Key idea:} Use the transmon qubit coupled to an SRF cavity
as a quantum sensor for the $\psi$-field, exploiting the extreme
sensitivity of the qubit's coherence to tiny perturbations.

\textbf{Mechanism:}
\begin{enumerate}
  \item The DFD $\psi$-field couples to the electromagnetic field
    inside the SRF cavity through the refractive index
    $n = e^\psi \approx 1 + \psi$.
  \item This induces a tiny frequency shift
    $\delta\omega/\omega \sim \psi \sim 10^{-9}$ on the cavity
    resonance.
  \item The transmon qubit, dispersively coupled to the cavity,
    experiences this frequency shift as a Stark shift on its
    transition frequency.
  \item With the cavity $Q \sim 10^{10}$ and qubit coherence
    $T_2 \sim 20$ ms, the accumulated phase over one coherence
    time is:
    \begin{equation}
      \Delta\phi = \delta\omega \cdot T_2
      \sim 2\pi \times 6\text{ GHz} \times 10^{-9} \times 20\text{ ms}
      \sim 0.75 \text{ rad}.
    \end{equation}
  \item This is a DETECTABLE phase shift using standard Ramsey
    interferometry.
\end{enumerate}

\textbf{Predicted signal:}
\begin{equation}
  \Delta\phi_\psi = 2\pi\,\nu_{\rm cavity}\,\psi_{\rm local}\,T_2
  \approx 0.75 \text{ rad} \times \left(\frac{\psi}{10^{-9}}\right)
  \left(\frac{T_2}{20\text{ ms}}\right).
\end{equation}

\textbf{Background discrimination:}
\begin{itemize}
  \item Modulate by rotating the cavity orientation relative to
    Earth's gravitational field (the $\psi$-gradient points radially).
  \item Compare two cavities at different heights (differential
    $\psi$-measurement, analogous to gravitational redshift).
  \item The DFD-specific signature: the phase shift should correlate
    with the Newtonian potential gradient, not with any EM background.
\end{itemize}

\textbf{Feasibility assessment:}
\begin{itemize}
  \item Required: two identical SRF cavity-qubit systems separated
    by $\Delta h \sim 1$ m.
  \item Expected $\delta\psi \sim g\Delta h/c^2 \sim 10^{-16}$.
  \item This is $10^7$ below the single-cavity sensitivity.
  \item \textbf{Challenge:} Differential measurement requires
    extraordinary common-mode rejection. Current technology is
    $\sim 3$--4 orders of magnitude short.
\end{itemize}

\textbf{Near-term approach (achievable with current SQMS technology):}
\begin{itemize}
  \item Instead of detecting the $\psi$-field directly, use the
    Q-ratio anomaly (P11's primary prediction).
  \item Measure the ratio of quality factors between BCS-paired and
    unpaired electron states.
  \item DFD predicts: $Q_{\rm ratio} = 0.52 \pm 0.08$ vs.\ BCS
    prediction $3.60 \pm 0.10$.
  \item The transmon qubit provides a more precise measurement of
    the cavity $Q$ than traditional RF techniques, through photon
    lifetime measurements.
\end{itemize}

\textbf{Grade: B --- Theoretically sound but experimentally
challenging for direct $\psi$ detection. Near-term Q-ratio
measurement is feasible.}

\end{tcolorbox}

\subsubsection{SQMS Q-Ratio: Status of Prediction}

The P11 prediction remains:
\begin{itemize}
  \item DFD: $Q_{\rm ratio} = 0.52 \pm 0.08$.
  \item BCS standard: $Q_{\rm ratio} = 3.60 \pm 0.10$.
  \item Separation: $> 30\sigma$.
  \item SQMS bound on DFD: $|\lambda - 1| < 2.0\ee{-13}$.
\end{itemize}

The new transmon-qubit measurements provide additional precision
on the cavity $Q$, but the Q-ratio test requires comparing paired
vs.\ unpaired states, which is a different experimental configuration.

\subsubsection{Disformal Framework Verification for P11}

From the i25 disformal $\Sigma(y)$ resolution (Cosmo update, Agent 4):
the disformal factor $D(X)$ contributes only at relativistic order.
For the SRF cavity (non-relativistic EM fields, $v/c \sim 0$):
\begin{equation}
  D\text{-contribution} \sim D(X)\,\partial_0\phi\,\partial_0\phi
  \sim D \cdot c^2 \cdot \psi^2 \sim 10^{-18},
\end{equation}
which is negligible compared to the conformal contribution
$C \sim 1 + O(\psi) \sim 1 + 10^{-9}$.

\textbf{Result:} The P11 Q-ratio prediction is UNCHANGED in the
disformal framework. The disformal factor is irrelevant for
non-relativistic laboratory physics.

\subsection{New Literature Reviewed}

\begin{tcolorbox}[colback=green!5!white, colframe=green!50!black,
  title=\textbf{Agent 8: NEW LITERATURE REVIEWED}]

\begin{enumerate}

\item \textbf{SQMS/Fermilab (2025), ``Ultracoherent superconducting
  cavity-based multiqudit platform with error-resilient control,''
  arXiv:2506.03286.}
  Two-cell SRF cavity with transmon qubit. $Q \sim 10^{10}$,
  coherence lifetime $> 20$ ms. Universal control of two normal modes.
  \emph{Directly relevant:} provides the hardware platform for the
  P11 $\psi$-probe concept.

\item \textbf{SQMS/Fermilab (2025), ``Near-term Application
  Engineering Challenges in Emerging Superconducting Qudit Processors,''
  arXiv:2506.05608.}
  Technical challenges for scaling SRF cavity-qubit systems.
  Identifies inverse Purcell decay, dielectric loss, and seam loss
  as the main limitations. All are materials/engineering challenges,
  not $\psi$-related.

\item \textbf{SQMS Center (2025--2026), ``Surface encapsulation of
  transmon qubits.''} Air-exposure-robust oxide growth suppresses
  TLS density, enabling multifold $Q$ increase and order-of-magnitude
  $T_1$ improvement. Critical for P11 sensitivity.

\item \textbf{Physics World (2025), ``Superconducting innovation:
  SQMS shapes up for scalable success in quantum computing.''}
  Overview of SQMS achievements and roadmap. Highlights the
  longest-lived multimode QPU ($> 20$ ms coherence).

\item \textbf{Analytical quantum full-wave analysis of few-photon
  transport through a superconducting cavity qubit (2026),
  arXiv:2603.03015.}
  Theoretical framework for photon transport in cavity-qubit systems.
  Relevant for understanding the quantum measurement protocols
  needed for the P11 $\psi$-probe.

\end{enumerate}
\end{tcolorbox}


%=============================================================================
\section{Appendix: Complete New Literature Table (Patent Team)}
%=============================================================================

\begin{longtable}{p{0.08\textwidth}p{0.50\textwidth}p{0.15\textwidth}p{0.15\textwidth}}
\toprule
\textbf{Agent} & \textbf{Paper} & \textbf{arXiv/Ref} & \textbf{Grade} \\
\midrule
\endhead
6 & Park et al., LMCP cloaking & Industry (2025) & A (dead) \\
6 & IMDEA, mechanical cloaking & Nat.\ Comm.\ (2025) & A (dead) \\
6 & AI acoustic metasurfaces & JSV (2025) & A (dead) \\
6 & SQMS surface encapsulation & SQMS (2025) & C \\
\midrule
7 & Saha et al., multi-messenger lensing & 2503.19973 & B \\
7 & Ferraro et al., GLPV GW signatures & 2509.05027 & B \\
7 & Atom interferometry Einstein Elevator & Nat.\ Comm.\ (2025) & B \\
7 & Cold Atom Lab ISS pathfinder & Nat.\ Comm.\ (2024) & C \\
7 & ET new baseline design & ET Collab.\ (2025) & B \\
\midrule
8 & SQMS multiqudit platform & 2506.03286 & B \\
8 & SQMS qudit engineering challenges & 2506.05608 & C \\
8 & SQMS surface encapsulation & SQMS (2025) & B \\
8 & Quantum photon transport & 2603.03015 & C \\
\bottomrule
\end{longtable}


\end{document}
